% ============================================================================
% Documentación: Vibration Damping Optimization (VDO) para CARP
% Generado/actualizado: 2026-07-10. Fuente: metacarp/vibration_damping.py,
% scripts/run_vdo_automatico.py.
% ============================================================================
\section{Vibration Damping Optimization (VDO)}
\label{sec:mh-vdo}

\subsection{Concepto algorítmico}
El VDO (Mehdizadeh y Tavakkoli-Moghaddam, 2008/2009) es una metaheurística de
trayectoria inspirada en el \textbf{amortiguamiento de la amplitud
vibratoria} en sistemas mecánicos: un oscilador que arranca con gran amplitud
$A_0$ pierde energía por fricción y su amplitud decae exponencialmente. Su
estructura es prima hermana del SA, pero el parámetro de control no es la
temperatura sino la \textbf{amplitud} $A$, y la regla de aceptación no es
Metropolis sino la CDF de una distribución de Rayleigh.

\paragraph{Regla de aceptación (Rayleigh).} Con $\Delta = c_{vecino} -
c_{actual}$: si $\Delta \le 0$ se acepta siempre; si $\Delta > 0$ se acepta
con probabilidad
\[
  p \;=\; 1 - \exp\!\left(-\frac{A^2}{2\sigma^2}\right),
\]
que es la CDF Rayleigh de escala $\sigma$ evaluada en $A$. Cuando $A \gg
\sigma$ (inicio), $p \approx 1$ (exploración agresiva); cuando $A \ll \sigma$
(final), $p \approx 0$ (explotación local). A diferencia del Metropolis del
SA, la regla \textbf{no depende de la magnitud de $\Delta$}: mientras el
sistema vibra con amplitud $A$, todo empeoramiento se arriesga con la misma
probabilidad. Esta es la diferencia conceptual clave entre VDO y SA.

\paragraph{Ley de amortiguamiento.} Tras cada nivel de $L$ iteraciones:
\[
  A_{t} \;=\; A_0 \cdot \exp\!\left(-\frac{\gamma\, t}{2}\right),
\]
la ecuación del oscilador amortiguado clásico evaluada en niveles discretos
$t = 0, 1, 2, \dots$ Siguiendo la letra del artículo original, el decaimiento
se recalcula \emph{siempre desde} $A_0$ (no acumulando un factor sobre el
nivel anterior, como hace el SA con $\alpha \times T$). La búsqueda se
detiene cuando $A$ cae bajo \texttt{umbral\_amplitud\_minima} o se agota el
presupuesto.

\subsection{Parámetros instance-aware}
Sea $n$ el número de arcos requeridos y $d_{\max}$ la distancia máxima de la
matriz Dijkstra:

\begin{table}[htbp]
\centering
\caption{Parámetros del VDO y defaults instance-aware.}
\begin{tabular}{lll}
\toprule
Parámetro & Símbolo & Default \\
\midrule
Amplitud inicial & $A_0$ & $20\, d_{\max}/n$ \\
Umbral de amplitud mínima & $A_{\min}$ & $20\, d_{\max}/n^2$ \\
Escala de Rayleigh & $\sigma$ & $A_0/2$ \\
Coef.\ de amortiguamiento & $\gamma$ & 0.05 \\
Iteraciones por nivel & $L$ & $n^2$ \\
Prob.\ inter-ruta (factible) & $p_{inter}$ & 0.6 \\
\bottomrule
\end{tabular}
\end{table}

$A_0$ y $A_{\min}$ replican la escala de $T_{init}$/$T_{\min}$ del SA para
que ambos regímenes exploratorios sean comparables; $\sigma = A_0/2$
garantiza probabilidad inicial de aceptación $1 - e^{-2} \approx 0.865$;
$\gamma$ es adimensional (constante del sistema oscilante) y no requiere
escalado por instancia. La implementación comparte con el resto del proyecto
los 9 operadores de vecindario, el objetivo penalizado con
$\lambda = \max(10 \cdot \text{mediana}(d), 10)$ y el sesgo inter/intra con
piso $0.8$ bajo violación de capacidad. Deliberadamente \emph{no} incluye los
mecanismos experimentales del SA (reheat, kick): es el VDO canónico.

\subsection{Búsqueda de malla (grid search)}
\begin{itemize}
  \item \textbf{Grilla} (\texttt{run\_vdo\_automatico.py}):
    $\gamma \in \{0.02, 0.05, 0.1, 0.2\}$ (4) $\times$
    \texttt{factor\_sigma} $\in \{0.25, 0.5, 1.0, 2.0\}$ (4) $\times$
    $p_{inter} \in \{0.4, 0.5, 0.6, 0.7, 0.8\}$ (5)
    $= 80$ combinaciones $\times$ 23 instancias $\times$ 5 repeticiones
    $= 9{,}200$ corridas. \texttt{factor\_sigma} escala $\sigma$
    proporcionalmente al $A_0$ efectivo de cada instancia (misma filosofía
    que \texttt{factor\_pasos} del CS), de modo que la grilla es comparable
    entre instancias de distinto tamaño; $A_0$ y $L$ se dejan en sus valores
    instance-aware.
  \item \textbf{Protocolo}: idéntico al resto de la campaña (semillas
    deterministas por tarea, \texttt{ProcessPoolExecutor}, CSV parciales
    consolidados por instancia, sin columnas \texttt{id\_corrida} /
    \texttt{config\_id}).
  \item \textbf{Estado}: el smoke-test inicial (julio 2026) igualó al SA en
    \texttt{gdb19} (costo 63 a presupuesto canónico, factibilidad completa
    c1--c5); la ejecución completa de la grilla y su calibración 2-knob están
    pendientes de re-ejecución, por lo que VDO aún no aparece en las tablas
    comparativas de la campaña costo-corregido.
\end{itemize}
